\section{Structural Features}
% ════════════════════════════════════════════════════════════

\textbf{Normal fault (primary):}
A well-exposed, east-dipping normal fault was traced for 1.4\,km along the
northern escarpment face.
Fault plane attitude: 082\textdegree/62\textdegree\ (strike/dip, right-hand rule).
Slickenlines rake 82\textdegree\ (near down-dip), indicating predominantly dip-slip motion.
Maximum measured stratigraphic offset: \textbf{42\,m}; gouge zone 15--30\,cm thick;
breccia zone up to 0.8\,m wide.
The hanging-wall block is rotated 8--12\textdegree\ relative to the footwall.

\textbf{Conjugate fault set:}
A secondary NW-trending normal fault set (mean attitude 322\textdegree/45\textdegree) produces
graben blocks 80--200\,m wide in the central section. Offsets 2--15\,m.

\textbf{Bedding attitudes:}
Mean attitude of Units I--V: 094\textdegree/08\textdegree\ (shallow eastward dip, consistent with
half-graben rotation). Units VI--VII: sub-horizontal (002\textdegree/03\textdegree), unconformably
overlying tilted lower units; angular unconformity well-exposed at Station~6.

\textbf{Joints:}
Two prominent joint sets: Set~A (N\,10\textdegree{}E, near-vertical); Set~B (N\,75\textdegree{}W, vertical).
Spacing 0.3--1.5\,m; open fractures filled with calcite veins up to 8\,mm.

% ════════════════════════════════════════════════════════════
